\section{Deterministic runtime engine}
\label{sec:runtime}

The canonical run loop is \code{Sim.Runtime.Engine} (\code{src/sim/Runtime/Engine.jl}). It advances an integer-microsecond timeline, calls discrete-time \emph{sources} at event boundaries, and integrates the continuous-time plant between boundaries.

\subsection{Timebase and hybrid event timeline}
\label{sec:runtime-timebase}

All scheduling uses integer microseconds (\code{UInt64}). The helper \code{dt\_to\_us(dt\_s)} (\code{src/sim/Runtime/Timeline.jl}) converts a \code{Float64} interval in seconds to microseconds with an exactness requirement:
\begin{equation}
  \left|dt\_s - 10^{-6}\,\mathrm{round}(10^{6} dt\_s)\right| \le 10^{-15}.
\end{equation}
If the conversion is not exact (i.e. the requested interval is not an integer multiple of \SI{1}{\micro\second}), the code throws.

\paragraph{Timeline axes.}
The runtime constructs several named time axes, each a sorted vector of microseconds (\code{src/sim/Runtime/Timeline.jl}): autopilot ticks, wind ticks, logging ticks, scenario \emph{AtTime} boundaries, and an optional ``physics'' axis. The global event axis is the sorted unique union of all axis times, including both endpoints. The engine always integrates on this union axis: any subsystem boundary becomes an integration boundary.

\begin{figure}[!htbp]
  \centering
  \begin{tikzpicture}[x=0.55cm,y=0.75cm, font=\small]
    \def\xmax{12}

    % Axis baselines
    \foreach \y/\name in {3/{autopilot axis},2/{wind axis},1/{log axis},0/{event axis (union)}} {
      \draw[-{Latex[length=2mm]}] (0,\y) -- (\xmax+0.8,\y);
      \node[anchor=east] at (-0.2,\y) {\name};
    }

    % Example tick sets (arbitrary units; the real code uses UInt64 microseconds).
    \foreach \t in {0,4,8,12} {\draw (\t,3+0.12) -- (\t,3-0.12);}
    \foreach \t in {0,3,6,9,12} {\draw (\t,2+0.12) -- (\t,2-0.12);}
    \foreach \t in {0,10,12} {\draw (\t,1+0.12) -- (\t,1-0.12);}
    \foreach \t in {0,3,4,5,6,8,9,10,12} {\draw (\t,0+0.12) -- (\t,0-0.12);}

    \node[anchor=west] at (0,-0.7) {\footnotesize Example only; in the implementation, all tick times are integer microseconds.};
  \end{tikzpicture}
  \caption{Multi-rate scheduling: each subsystem owns a time axis, and the engine integrates on their union (the event axis).}
  \label{fig:timeline-union}
\end{figure}

\subsection{TOML configuration: timeline axes}
\label{sec:runtime-toml}

The timeline parameters described above are configured in the \code{[timeline]} block of the aircraft spec (\code{TimelineSpec} in \code{src/sim/Aircraft/Spec.jl}, parsed by \code{src/sim/Aircraft/TOMLIO.jl}):

\begin{lstlisting}[language=toml]
[timeline]
t_end_s = 20.0
dt_autopilot_s = 0.004
dt_wind_s      = 0.001
dt_log_s       = 0.01
dt_phys_s      = 0.001  # optional; omit (or set to null) to disable extra physics boundaries
\end{lstlisting}

\paragraph{Parameter semantics (what the code does).}
\begin{itemize}[leftmargin=*,itemsep=2pt]
  \item \code{t\_end\_s}: inclusive simulation horizon. The engine always starts at \(t_0=0\) and forces both \(t_0\) and \(t_{end}\) into the union event axis.
  \item \code{dt\_autopilot\_s}: cadence of the autopilot axis (PX4 stepping, estimator publication, plant discontinuities).
  \item \code{dt\_wind\_s}: cadence of the wind axis. Wind is updated only on these ticks and held between them.
  \item \code{dt\_log\_s}: cadence of logging ticks. The implementation additionally appends a final log sample at \(t_{end}\) even when \(dt_{log}\) does not evenly divide the horizon (\code{build\_timeline}).
  \item \code{dt\_phys\_s}: optional additional periodic boundaries. When omitted, the ``physics'' axis degenerates to \(\{t_0,t_{end}\}\) and does not add boundaries beyond the other axes.
\end{itemize}

\paragraph{Mapping to the engine axes.}
Each interval is converted to integer microseconds by \code{dt\_to\_us} (\cref{sec:runtime-timebase}), giving
\begin{equation}
  \Delta t_{ap} \leftrightarrow \texttt{dt\_autopilot\_s},\quad
  \Delta t_{wind} \leftrightarrow \texttt{dt\_wind\_s},\quad
  \Delta t_{log} \leftrightarrow \texttt{dt\_log\_s},\quad
  \Delta t_{phys} \leftrightarrow \texttt{dt\_phys\_s}.
\end{equation}
Scenario \emph{AtTime} boundaries are \emph{not} configured in \code{[timeline]}; they are discovered from the scenario source via the \code{event\_times\_us} protocol hook and injected into the union axis by \code{build\_timeline\_for\_run}.

\subsection{Boundary protocol and sample-and-hold semantics}
\label{sec:runtime-boundary}

At each event time $t_k$ the engine executes boundary stages in a fixed canonical order (\code{src/sim/Runtime/BoundaryProtocol.jl}, applied by \code{process\_events\_at!} in \code{src/sim/Runtime/Engine.jl}):
\begin{enumerate}[leftmargin=*,itemsep=2pt]
  \item \textbf{Scenario}: update high-level commands (arm/mission/RTL), landed flag, faults, and optional wind disturbances.
  \item \textbf{Wind}: if the wind axis is due, step the wind source and publish the held \code{bus.wind\_ned} sample.
  \item \textbf{Derived outputs}: if the plant implements \code{plant\_outputs}, evaluate algebraic outputs (notably battery telemetry) and publish them into the bus.
  \item \textbf{Estimator}: if the autopilot axis is due, publish \code{bus.est} (truth injection + noise/bias/delay).
  \item \textbf{Telemetry}: optional deterministic inspection hooks (autopilot-axis only).
  \item \textbf{Autopilot}: if the autopilot axis is due, step PX4 (or replay) and publish \code{bus.cmd}.
  \item \textbf{Plant discontinuities}: if the autopilot axis is due, apply boundary-time plant updates (notably snapping \emph{direct} actuators via \code{plant\_on\_autopilot\_tick}).
  \item \textbf{Logging}: if the log axis is due, record a pre-step snapshot and emit structured logs.
\end{enumerate}

\paragraph{Why this order?}
The boundary protocol is part of the simulator semantics, not an implementation detail. At each boundary \(t_k\) the engine applies a discrete transition
\[
  (\mathrm{bus}_{k^-}, x_k) \;\xrightarrow{\;\text{boundary stages}\;}\; (\mathrm{bus}_{k^+}, x_k^+),
\]
where \(x_k\) is the continuous plant state at the boundary and \(x_k^+\) includes any boundary-time discontinuities (e.g.\ actuator snaps). The held plant input over the next interval is then defined by
\[
  u_k \coloneqq \mathrm{PlantInput}(\mathrm{bus}_{k^+}),\qquad
  u(t)=u_k\ \text{for}\ t\in[t_k,t_{k+1}).
\]
Reordering stages changes which signals are held over which intervals and therefore changes the meaning of the simulation.

The chosen ordering enforces the following invariants at PX4 ticks:
\begin{itemize}[leftmargin=*,itemsep=2pt]
  \item \textbf{No stale estimator inputs:} when the autopilot axis is due, PX4 consumes an estimator packet constructed from truth at the same boundary time \(t_k\).
  \item \textbf{No stale telemetry:} algebraic telemetry (e.g.\ battery status) is computed from \((t_k,x_k,u_k)\) and published before stepping PX4, so uORB injection reflects the current plant state.
  \item \textbf{Command-to-physics causality:} PX4 is stepped exactly once at \(t_k\), producing the unique command \(u_k\) that is held over \([t_k,t_{k+1})\).
  \item \textbf{Snap-after-command:} boundary-time discontinuities are applied after \code{bus.cmd} is published, so snapped actuator state is a deterministic function of the newly produced command.
  \item \textbf{Logs correspond to integration inputs:} logging occurs last so a snapshot at \(t_k\) corresponds to the same bus packet that seeds integration over \([t_k,t_{k+1})\).
\end{itemize}

Changing this ordering is therefore a semantic change (not ``just refactoring''): for example, stepping PX4 before the estimator would make PX4 consume a stale estimate; logging before PX4 would decouple recorded commands from the plant inputs that followed.

\paragraph{Scenario cadence.}
The scenario stage is currently executed at \emph{every} boundary on the union event axis, not only on the scenario axis. This makes hybrid \code{When(...)}-style logic responsive at the earliest possible boundary while preserving deterministic ordering (see the note in \code{src/sim/Runtime/BoundaryProtocol.jl} and the implementation in \code{process\_events\_at!}).

\paragraph{Wind disturbance consistency.}
If the scenario updates an additive wind disturbance at a boundary where the wind axis is \emph{not} due, the engine applies the disturbance delta immediately to the held wind sample so the piecewise-constant wind input remains consistent across axes.

\paragraph{Sample-and-hold.}\label{sec:runtime-sample-hold}
Between event times, all bus values are treated as sample-and-hold. The plant is integrated over $[t_k, t_{k+1})$ with a single input packet
\begin{equation}
  u_k \coloneqq \{\text{actuator command},\; \text{wind},\; \text{fault state}\},
\end{equation}
constructed from the bus at $t_k$ (\code{PlantInput} in \code{src/sim/Plant.jl}).


\subsection{Integration stepping and post-step projection}

Let $x(t)$ be the continuous plant state and $f(t, x, u)$ the plant RHS functor. At each interval the engine calls
\begin{equation}
  x_{k+1} \gets \mathrm{Integrate}(f,\, t_k,\, x_k,\, u_k,\, \Delta t_k)
\end{equation}
where $\Delta t_k = (t_{k+1}-t_k)\,10^{-6}$ seconds, and the integrator is selected by the aircraft spec (\code{src/sim/Integrators.jl}, \code{src/sim/Aircraft/Build.jl}).

\paragraph{Hybrid interval integration (contact-aware plants).}
\label{sec:runtime-hybrid-integrators}
Most plants can be advanced by calling \code{Integrators.step\_integrator} once per union-axis interval. Ground contact, however, is inherently hybrid: touchdown and liftoff introduce discontinuities that are not representable as a smooth ODE without either (i) stiff penalty forces or (ii) explicit event handling. The runtime therefore supports an optional protocol hook
\begin{quote}\small
\code{plant\_integrate\_interval(model, integrator, t0\_us, x0, u, dt\_us)}
\end{quote}
(\code{src/sim/PlantInterface.jl}). When implemented, the engine delegates the \emph{entire} interval advance to this method.

\paragraph{Free flight vs.\ contact-constrained stepping (current ground-contact implementation).}
For the coupled multirotor plant, \code{plant\_integrate\_interval} is implemented when \code{PlantSpec.contact} is \code{FlatGroundConstraintContact}. The interval integrator acts as a deterministic regime switcher:
\begin{itemize}[leftmargin=*,itemsep=2pt]
  \item \textbf{Airborne segments (smooth dynamics):} the plant is advanced under the \emph{free} dynamics \(f_{\mathrm{free}}\) (the same model but with \code{NoContact()}) using the configured integrator (RK4/RK23/RK45, etc.). Near the ground, the next airborne chunk can be capped to \(O(\text{time-to-impact})\) so that a descending crossing of the guard is bracketed reliably.
  \item \textbf{Impact events (discontinuous map):} if a lookahead airborne chunk brackets a descending crossing of \(z=0\), the TOI is localized to integer microseconds by deterministic bisection and an impact map is applied at that internal time (position projection + velocity reset with restitution and optional tangential friction). Multiple impacts can be handled inside one outer interval, but the code caps the count to a fixed maximum to avoid pathological bounce storms.
  \item \textbf{Grounded segments (constraint/impulse stepping):} while the state is a ground candidate, the plant advances in fixed substeps of \code{h\_contact\_us}. Each substep uses a single evaluation of \(f_{\mathrm{free}}\) (explicit Euler) followed by a closed-form, 1-contact impulse solve (normal + Coulomb friction) and a position-only stabilization on \(z\). In this grounded regime the configured integrator is \emph{not} used; the time-stepping loop is deliberately fixed-step and deterministic.
\end{itemize}
The full contact math and the exact update equations are documented in \cref{sec:contacts-hybrid}.

\paragraph{Interval metadata and logging.}
Because impacts can occur at internal TOI times \(\tau\in(t_k,t_{k+1})\), the interval integrator may optionally return per-interval metadata (impact count, maximum \(\Delta v\), and its timestamp). The engine accumulates this metadata into latched bus signals for logging (\code{bus.impact\_dv\_ned}, \code{bus.impact\_time\_us}, \code{bus.impact\_count}) and also computes a simple ``equivalent acceleration'' estimate at the autopilot tick rate (\code{bus.impact\_accel\_est\_ned}); see \code{Runtime.Engine.\_accumulate\_interval\_impacts!}.

Some state variables have hard physical bounds (e.g. rotor speed $\omega \ge 0$, battery SOC in $[0,1]$). The engine therefore supports an optional post-step projection hook \code{plant\_project(model, x)} that clamps or projects state deterministically after each accepted interval. The coupled multirotor plant implements this projection (\code{src/sim/PlantModels/CoupledMultirotor.jl}).

\paragraph{Record $\to$ replay (high level).}
The same engine can run in a mode that records boundary-aligned bus/plant streams, and a replay mode that replaces ``live'' sources with trace samplers. This converts a closed-loop PX4 run into a deterministic open-loop plant replay for integrator/model comparisons.

\begin{figure}[t]
  \centering
  \begin{tikzpicture}[
      node distance=6mm and 10mm,
      font=\small,
      box/.style={draw, align=center, minimum width=40mm, inner ysep=5pt, inner xsep=6pt},
      arr/.style={-Latex, line width=0.4pt}
    ]
    \node[box] (b) {Boundary at $t_k$\\(sources + logging)};
    \node[box, below=of b] (i) {Integrate plant\\$[t_k, t_{k+1})$ with held $u_k$};
    \node[box, below=of i] (p) {Optional projection\\$x \leftarrow \mathrm{plant\_project}(x)$};
    \draw[arr] (b) -- (i);
    \draw[arr] (i) -- (p);
  \end{tikzpicture}
  \caption{Engine rhythm: boundary processing on the union event axis, then continuous plant integration over the next interval under sample-and-hold inputs.}
  \label{fig:engine-rhythm}
\end{figure}
